\documentclass[11pt,a4paper]{article}

\usepackage[T1]{fontenc}
\usepackage[utf8]{inputenc}
\usepackage[english]{babel}
\usepackage{lmodern}
\usepackage{microtype}
\usepackage[a4paper,margin=2.2cm]{geometry}
\usepackage{amsmath,amssymb}
\usepackage{graphicx}
\usepackage{booktabs}
\usepackage{tabularx}
\usepackage{array}
\usepackage{enumitem}
\usepackage{xcolor}
\usepackage{hyperref}
\usepackage[font=small,labelfont=bf]{caption}
\usepackage{float}
\usepackage[most]{tcolorbox}
\usepackage{natbib}
\usepackage{siunitx}

\definecolor{Accent}{HTML}{1F4E79}
\definecolor{AccentLight}{HTML}{EAF2F8}
\hypersetup{
    colorlinks=true,
    linkcolor=Accent,
    citecolor=Accent,
    urlcolor=Accent,
    pdftitle={XY-Invariant Neural Networks for Sparse Timepix Particle Classification},
    pdfauthor={OpenAI},
    pdfsubject={Technical design note},
    pdfkeywords={particle classification, Timepix, sparse hits, point clouds, rotation invariance}
}

\setlength{\parindent}{0pt}
\setlength{\parskip}{0.55em}
\setlist[itemize]{leftmargin=1.45em,itemsep=0.25em,topsep=0.25em}
\setlist[enumerate]{leftmargin=1.7em,itemsep=0.25em,topsep=0.25em}
\sisetup{group-separator={,},group-minimum-digits=4}

\newcolumntype{Y}{>{\raggedright\arraybackslash}X}
\newcolumntype{L}[1]{>{\raggedright\arraybackslash}p{#1}}
\newcommand{\code}[1]{\texttt{#1}}
\newcommand{\thetaxy}{\theta_{xy}}
\newcommand{\phit}{\phi_t}

\begin{document}

{\color{Accent}\rule{\textwidth}{1.5pt}}
\vspace{0.6em}

{\LARGE\bfseries XY-Invariant Neural Networks for Sparse Timepix Particle Classification\par}
\vspace{0.35em}
{\large Updated technical design for turning the repository into a runnable NN classification testbed\par}

\vspace{0.6em}
\begin{tabularx}{\textwidth}{@{}Y Y@{}}
\textbf{Primary design source:} \code{xy\_invariant\_particle\_report\_project.zip} &
\textbf{Updated:} \today \\
\end{tabularx}
\vspace{0.8em}

\begin{abstract}
This report replaces the earlier report-first recommendation with an implementation-oriented design contract. Candidate particles should be represented as variable-size point sets
\[
P=\{(x_i,y_i,t_i,e_i)\}_{i=1}^{M},
\]
where $x,y$ are detector-plane coordinates, $t$ is a time-like coordinate derived from ToA/FToA or frame metadata, and $e$ is an energy-like feature such as $\log(1+\mathrm{ToT})$ until calibrated energy is available. The central correction is selective invariance: detector-plane rotation $\thetaxy$ must be computed and stored, but it must not define a class. A horizontal, vertical, or diagonal copy of the same track-like morphology should map to the same class embedding. Time pitch $\phit$, shape, energy density, scale, and total energy remain classification-relevant.

The recommended baseline is a deterministic geometry front-end plus a compact canonical PointNet/DeepSets model called \textbf{XYInvariantParticleNet}. DBSCAN remains useful as an interpretable candidate generator, but clustering/classification after the encoder must operate only on $z_{class}$, never on a saved descriptor that includes $\thetaxy$.
\end{abstract}

\begin{tcolorbox}[colback=AccentLight,colframe=Accent,title=\textbf{Main Design Contract}]
\textbf{Same class:} same morphology, similar energy-density profile, comparable scale, and similar time pitch, under any detector-plane rotation.\\
\textbf{Stored metadata only:} $\thetaxy$ and its reliability $q_\theta$ are saved for downstream physics, visualization, and direction statistics. They are excluded from class distance, HDBSCAN, DEC, VaDE, and classifier logits.
\end{tcolorbox}

\section{Data Inventory}

The repository now has two data layers. The small tracked example is \code{data/matrix\_dump\_0001.txt}, a sparse JSON-style dump containing 198 matrices across 40 \code{sample} blocks. The large raw archive \code{raw\_data.zip} is extracted locally under \code{local\_data/raw/} and is intentionally gitignored.

The raw archive contains tab-separated \code{.t3pa} hit tables paired with \code{.t3pa.info} sidecars. Each raw row has columns \code{Index}, \code{Matrix Index}, \code{ToA}, \code{ToT}, \code{FToA}, and \code{Overflow}. The current parser decodes \code{Matrix Index} as row-major coordinates:
\[
x = \mathrm{index}\bmod 256,\qquad y=\lfloor \mathrm{index}/256 \rfloor.
\]
The sidecars add acquisition metadata such as chipboard ID, high voltage, acquisition duration, threshold, Pixet version, and start time.

\begin{table}[H]
\centering
\caption{Raw archive orientation generated by the repository indexer.}
\label{tab:raw-index}
\begin{tabularx}{\textwidth}{@{}lrl@{}}
\toprule
\textbf{Item} & \textbf{Count / size} & \textbf{Notes} \\
\midrule
Tracked example dump & 198 matrices & Sparse $256\times256$ calibrated-energy matrices \\
Raw files & 711 \code{.t3pa} + 711 sidecars & Six run folders in \code{local\_data/raw/} \\
Uncompressed raw bytes & \SI{1.426}{GB} & Large data remains outside git \\
Raw hit rows & 44,725,206 & Counted directly from \code{.t3pa} rows \\
Index artifacts & CSV + Markdown & \code{data/raw\_data\_index.csv}, \code{docs/raw-data-index.md} \\
\bottomrule
\end{tabularx}
\end{table}

\section{Target Product}

The measured input is not yet a particle class. It is a set of detector hits. The useful reconstruction chain should keep these layers separate:
\begin{enumerate}
    \item hit filtering or denoising;
    \item clustering or instance segmentation into particle candidates;
    \item candidate-level morphology/energy classification;
    \item optional association across time, views, or chipboards;
    \item run-level summaries and physics validation.
\end{enumerate}

The first practical deliverable is therefore:
\[
\text{raw hits} \rightarrow \text{candidate clustering} \rightarrow \text{XY-invariant candidate descriptor and class confidence}.
\]
The returned candidate record should include at least source path or \code{sample}/\code{set\_index}, candidate ID, hit list or mask, total energy/ToT, size, eccentricity or PCA ratios, $z_{class}$, class confidence or cluster ID, $\thetaxy$, $q_\theta$, and $\phit$.

\section{Selective XY Invariance}

Detector-plane azimuth is a nuisance factor for morphology classification. A straight object at $0^\circ$, $45^\circ$, and $90^\circ$ should not become three classes. Conversely, the time axis is not another spatial axis. The pitch between the spatial trajectory and time can encode a physically meaningful difference and should remain available to the class descriptor.

\begin{table}[H]
\centering
\caption{Which factors may affect classification or clustering.}
\label{tab:factors}
\begin{tabularx}{\textwidth}{@{}L{0.30\textwidth}L{0.20\textwidth}Y@{}}
\toprule
\textbf{Factor} & \textbf{Use in class distance?} & \textbf{Reason} \\
\midrule
Shape latent $z_{shape}$ & Yes & Defines point-like, blob-like, track-like, curved, shower-like, or fragmented morphology \\
Energy-density latent $z_{density}$ & Yes & Separates shapes with different ToT/energy profiles \\
Scale $z_{scale}$ & Yes, configurable & Length, width, duration, and hit count may be class-relevant \\
Time pitch $z_{tilt}$ / $\phit$ & Yes & Time is physically special and should not be mixed with detector-plane rotation \\
Energy summary $z_{energy}$ & Yes & Needed for energetic taxonomy and downstream validation \\
Detector-plane azimuth $\thetaxy$ & \textbf{No} & Store and analyze, but never use to split morphology classes \\
Angle reliability $q_\theta$ & No & Metadata indicating whether $\thetaxy$ is meaningful \\
\bottomrule
\end{tabularx}
\end{table}

\section{Geometry Front-End}

For a candidate $P=\{(x_i,y_i,t_i,e_i)\}_{i=1}^{M}$, compute an energy-weighted centroid and covariance. Let $w_i=e_i+\epsilon$ and use a time scale $\lambda_t$ to place the time feature in a controlled numerical range:
\[
r_i=(x_i-\bar{x},\; y_i-\bar{y},\; \lambda_t(t_i-t_{min})).
\]
The weighted covariance eigenvectors provide two useful angles:
\[
\thetaxy=\operatorname{atan2}(v_y,v_x),\qquad
\phit=\operatorname{atan2}\left(v_t,\sqrt{v_x^2+v_y^2}\right).
\]
For compact or one-hit candidates, detector-plane angle can be meaningless. Store a reliability score:
\[
q_\theta=\frac{\lambda_{xy,1}-\lambda_{xy,2}}{\lambda_{xy,1}+\lambda_{xy,2}+\epsilon}.
\]
For the class branch, rotate only $x,y$ into a canonical frame:
\[
\begin{pmatrix}x'_i\\y'_i\end{pmatrix}
=R(-\thetaxy)
\begin{pmatrix}x_i-\bar{x}\\y_i-\bar{y}\end{pmatrix}.
\]
The original $\thetaxy$ is preserved in metadata, not in the clustering vector.

\section{XYInvariantParticleNet}

The first neural model should be deliberately modest: a PointNet/DeepSets encoder over canonicalized variable-size point sets \citep{qi2017pointnet,zaheer2017deepsets}. A stronger local model can later add EdgeConv or point-transformer blocks \citep{wang2019dgcnn,zhao2021pointtransformer,robles2025pointset}.

\begin{table}[H]
\centering
\caption{Reference model contract.}
\label{tab:model}
\begin{tabularx}{\textwidth}{@{}L{0.29\textwidth}Y@{}}
\toprule
\textbf{Component} & \textbf{Role} \\
\midrule
Input candidate & Variable-size set $[M,F]$ with canonicalized $x',y'$, relative time, and $\log(1+\mathrm{ToT})$ or calibrated energy \\
Geometry front-end & Computes centroid, scale, weighted PCA, $\thetaxy$, $q_\theta$, $\phit$, and energy-density profiles \\
Hit MLP & Shared per-hit encoder for any number of hits \\
Pooling & Max, mean, and attention pooling produce a fixed-length candidate vector \\
Class latent & $z_{class}=[z_{shape},z_{density},z_{scale},z_{tilt},z_{energy}]$; excludes $\thetaxy$ \\
Pose metadata head & Predicts or validates $\thetaxy$ and $q_\theta$ for saved metadata only \\
Clustering layer & HDBSCAN, DEC, IDEC, VaDE, GMM, or prototype clustering over $z_{class}$ only \citep{mcinnes2017hdbscan,xie2016dec,guo2017idec,jiang2017vade} \\
\bottomrule
\end{tabularx}
\end{table}

The implementation rule is:
\[
\boxed{D(P_a,P_b)=D(z^a_{class},z^b_{class}) \quad\text{and never}\quad D(z^a_{saved},z^b_{saved}).}
\]
If $\thetaxy$ leaks into $D$, a later clustering algorithm will split horizontal, vertical, and diagonal copies of the same physical morphology.

\section{Training Without Labels}

The repository should first support self-supervised or weakly supervised learning. The useful loss family is:
\[
\mathcal{L}=\mathcal{L}_{profile}
+\alpha\mathcal{L}_{xy\text{-}inv}
+\beta\mathcal{L}_{xy\text{-}pose}
+\gamma\mathcal{L}_{tilt}
+\delta\mathcal{L}_{density}.
\]

The core invariance loss rotates the detector plane and requires the class embedding to stay stable:
\[
\mathcal{L}_{xy\text{-}inv}=\left\|z_{class}(P)-z_{class}(R_{\Delta\theta}P)\right\|_2^2.
\]
The pose head can still be trained to shift by the known augmentation angle:
\[
\hat{\theta}(R_{\Delta\theta}P)\approx \hat{\theta}(P)+\Delta\theta.
\]
For undirected line-like candidates, use double-angle encoding $(\cos 2\theta,\sin 2\theta)$ so $\theta$ and $\theta+\pi$ represent the same axis.

Useful self-supervised targets include longitudinal/radial energy profiles, ToT quantiles, PCA eigenvalue ratios, hit count, total ToT, and time span. A masked point autoencoder or profile autoencoder can later be added, following the same sparse point-set logic \citep{young2025polarmae}.

\section{Baselines and Roadmap}

\textbf{Stage 0: raw-data orientation.} Extract the raw archive into \code{local\_data/raw/}, generate \code{docs/raw-data-index.md}, and verify row counts, metadata, and run folders.

\textbf{Stage 1: deterministic XY-invariant descriptor.} Compute weighted PCA, energy profiles, scale, energy summaries, and $\phit$. Cluster these descriptors with HDBSCAN or DBSCAN, excluding $\thetaxy$. This is the sanity check against which the NN must improve.

\textbf{Stage 2: canonical PointNet/DeepSets.} Train XYInvariantParticleNet on DBSCAN candidates or curated candidate windows using rotation-invariance, pose, and profile losses. Report rotation consistency and descriptor leakage.

\textbf{Stage 3: stronger point models.} Add EdgeConv, GravNet-like latent neighborhoods, or point transformers when local shape, branches, and overlapping deposits require more capacity \citep{qasim2019gravnet,qu2020particlenet,qu2022part}.

\textbf{Stage 4: learned clustering.} If simulation or hit-level truth becomes available, replace the DBSCAN candidate generator with object condensation or a related instance-segmentation model \citep{kieseler2020objectcondensation,qasim2021hgcal_oc,garcia2026hitpf}.

\section{Validation}

Accuracy alone is not sufficient. The required tests are:
\begin{itemize}
    \item \textbf{rotation consistency:} variance of $z_{class}(R_\alpha P)$ over random $\alpha$;
    \item \textbf{pose accuracy:} circular error of predicted or stored $\thetaxy$;
    \item \textbf{cluster invariance:} fraction of rotated copies assigned to the same cluster;
    \item \textbf{descriptor leakage:} correlation or mutual information between cluster ID and $\thetaxy$ should be low;
    \item \textbf{tilt sensitivity:} $\phit$ changes can separate groups when intended;
    \item \textbf{interpretability:} cluster medoids should differ by shape, energy density, scale, or time pitch, not by detector-plane azimuth.
\end{itemize}

\section{Conclusion}

The repository should treat raw Timepix rows and sparse matrix dumps as measured detector hits, not ordinary images. The immediate implementation target is a runnable baseline that can parse raw data, build candidate descriptors, run DBSCAN as a bootstrap, and train a compact point-set encoder. The governing design principle is selective invariance: compute and preserve detector-plane orientation, but do not let it define the morphology class. This makes the system debuggable today and keeps a clean path toward stronger graph, transformer, or object-condensation models when labels or simulation truth become available.

\bibliographystyle{unsrtnat}
\bibliography{particle_nn_report}

\end{document}
